\documentclass[11pt,a4paper]{article}
\usepackage{amsmath}
\usepackage{fontspec}
\usepackage{babel}
\usepackage[margin=2.5cm]{geometry}
\usepackage{enumitem}
\usepackage{xcolor}
\usepackage{hyperref}
\usepackage{mdframed}

\definecolor{coral}{HTML}{cc785c}
\definecolor{ink}{HTML}{141413}
\definecolor{muted}{HTML}{6c6a64}
\definecolor{lightbg}{HTML}{f5f4f0}

\hypersetup{colorlinks=true, linkcolor=coral, urlcolor=coral}

\babelprovide[import, onchar=ids fonts]{thai}
\babelprovide[import, onchar=ids fonts]{english}

\babelfont{rm}{Noto Sans}
\babelfont[thai]{rm}{Noto Sans Thai}
\babelfont{tt}{Noto Sans Mono}
\babelfont[thai]{tt}{Noto Sans Thai}

\directlua{
  local glyph_id = node.id("glyph")
  local penalty_id = node.id("penalty")

  local function is_thai(c)
    return c >= 0x0E01 and c <= 0x0E5B
  end

  local function is_thai_combining(c)
    return (c >= 0x0E31 and c <= 0x0E3A) or (c >= 0x0E47 and c <= 0x0E4E)
  end

  luatexbase.add_to_callback("pre_linebreak_filter", function(head)
    for n in node.traverse_id(glyph_id, head) do
      local prev = n.prev
      if prev and prev.id == glyph_id and is_thai(prev.char)
         and is_thai(n.char) and not is_thai_combining(n.char) then
        local p = node.new(penalty_id)
        p.penalty = 0
        node.insert_before(head, n, p)
      end
    end
    return head
  end, "thai_break")
}

\emergencystretch=1em

\newmdenv[
  backgroundcolor=lightbg,
  linecolor=coral,
  linewidth=2pt,
  topline=false, rightline=false, bottomline=false,
  innerleftmargin=12pt, innerrightmargin=8pt,
  innertopmargin=8pt, innerbottommargin=8pt,
]{highlight}

\begin{document}

\begin{center}
  {\LARGE \textbf{สรุปบทความ}}\\[6pt]
  {\color{muted}\small Paper Reading Note}
\end{center}

\noindent\rule{\textwidth}{1.5pt}

%% ─── Metadata ─────────────────────────────────────────────────
\vspace{12pt}
\noindent
\begin{tabular}{@{}ll@{}}
  \textbf{ชื่อบทความ}  & ชื่อเต็มของบทความ \\[4pt]
  \textbf{ผู้เขียน}    & Author1, Author2, Author3 \\[4pt]
  \textbf{ปี / สถานที่} & 2024 · Conference/Journal Name \\[4pt]
  \textbf{arXiv / DOI}  & arXiv:XXXX.XXXXX \\[4pt]
  \textbf{อ่านเมื่อ}   & \today \\
\end{tabular}

\vspace{12pt}
\noindent\rule{\textwidth}{0.4pt}

%% ─── Problem ──────────────────────────────────────────────────
\section*{ปัญหาที่แก้ไข}

อธิบายปัญหาหลักที่บทความนี้แก้ไข ทำไมปัญหานี้จึงสำคัญ
และทำไมวิธีเดิมยังไม่เพียงพอ

%% ─── Contribution ─────────────────────────────────────────────
\section*{ผลงานสำคัญ (Contribution)}

\begin{enumerate}[nosep, leftmargin=2em]
  \item Contribution ที่ 1
  \item Contribution ที่ 2
  \item Contribution ที่ 3
\end{enumerate}

%% ─── Method ───────────────────────────────────────────────────
\section*{วิธีการ}

อธิบายแนวทางหลักของบทความนี้ รวมถึงสถาปัตยกรรม
หรือกระบวนการสำคัญ

%% ─── Key Results ──────────────────────────────────────────────
\section*{ผลลัพธ์สำคัญ}

\begin{highlight}
\textbf{ตัวเลขสำคัญ:}
\begin{itemize}[nosep]
  \item Dataset X: +XX\% over baseline (metric)
  \item Dataset Y: +XX\% over baseline (metric)
\end{itemize}
\end{highlight}

%% ─── Strength / Weakness ──────────────────────────────────────
\section*{จุดแข็ง / จุดอ่อน}

\noindent\textbf{จุดแข็ง:}
\begin{itemize}[nosep, leftmargin=2em]
  \item จุดแข็งที่ 1
  \item จุดแข็งที่ 2
\end{itemize}

\vspace{6pt}
\noindent\textbf{จุดอ่อน / ข้อจำกัด:}
\begin{itemize}[nosep, leftmargin=2em]
  \item ข้อจำกัดที่ 1
  \item ข้อจำกัดที่ 2
\end{itemize}

%% ─── Relevance ────────────────────────────────────────────────
\section*{ความเกี่ยวข้องกับงานของเรา}

อธิบายว่าบทความนี้เกี่ยวข้องกับ thesis ของเราอย่างไร
จะนำแนวคิดใดไปใช้ และต้องระวังอะไร

%% ─── Notes ────────────────────────────────────────────────────
\section*{บันทึกเพิ่มเติม}

หมายเหตุ คำถาม หรือสิ่งที่ต้องติดตามเพิ่มเติม

\vspace{16pt}
\noindent\rule{\textwidth}{1pt}
\noindent{\color{muted}\small สรุปโดย Claude Sonnet 4.6 · \today}

\end{document}
